% !TeX root = ../main.tex
\section{Task and Simulation Setup}
\label{sec:setup}
\subsection{Physical Task and Success Criteria}
The MuJoCo environment \citep{todorov2012mujoco} contains a passive hinged
door, a fixed vertical handle, and a mobile manipulator with a lift joint,
shoulder, elbow, wrist, and parallel two-finger gripper. An ideal holonomic
base moves in $x$, $y$, and yaw, then remains at its selected station during
opening. The mobile controller uses a 6~N actuator limit per finger, a
10-second opening command, and a 2-second hold phase. The target opening is
$90^\circ$ in the selected direction.

Success requires bilateral grasp acquisition and, throughout the final
second, a door angle within $2^\circ$ of the target, bilateral contact, and
angular speed below $0.05$~rad/s. Unexpected collision, sustained slip, or
contact loss terminates an ordinary trial. Slip monitoring uses relative
motion and contact-point velocity; persistent loss of bilateral contact for
0.1~s raises a contact-loss event. These checks distinguish reaching the
handle, acquiring a grasp, and completing the opening.

\subsection{Observations and Candidate Actions}
The robot starts with calibrated RGB-D and a 19-dimensional state record.
One optional inspection pose provides a physically reached second view. P
chooses among four predefined stations, two on each side of the door, or
refusal. After station arrival, G uses a fresh observation to select one of
two approach offsets, $-15^\circ$ and $+22^\circ$, with a side-dependent yaw
conversion. Grasp height, IK branch, grip force, and timing are fixed in the
mobile policy.

Opening constraints combine the door/robot CAD prior with observed depth.
Four features describe each station: known base collision, observed
obstruction, sampled free-space coverage, and normalized observed clearance.
Observed obstruction can exclude a station before learned ranking. Unseen
space remains unknown; navigation and arm preparation use a separate
known-geometry service. Table~\ref{tab:information} specifies the information
available at each decision.

\begin{table}[H]
\caption{Information and decisions along one episode. Reached images become
available only after the corresponding physical movement.}
\label{tab:information}
\centering\small
\begin{tabularx}{\linewidth}{@{}lXX@{}}
\toprule
Stage & Available information & Decision or execution \\
\midrule
I & Initial RGB features, relative inspection goal, risk, initial depth features
  & Inspect or continue directly \\
P & Current visual features, relative station commands, current depth features
  & Four stations or refusal \\
G & Reached-station image and estimated handle geometry
  & Two approach offsets \\
Controller & Selected plan, joint state, known geometry, contact feedback
  & Navigate, prepare arm, grasp, open, hold \\
\bottomrule
\end{tabularx}
\end{table}
